% Generated draft: requires mathematical and source review.
\documentclass[11pt]{article}
\usepackage[margin=25mm]{geometry}
\usepackage{fontspec}
\setmainfont{Latin Modern Roman}
\usepackage{amsmath,amssymb,mathtools,unicode-math}
\setmathfont{Latin Modern Math}
\usepackage{enumitem,needspace,xurl,hyperref}
\setlength{\parindent}{0pt}
\setlength{\parskip}{5pt}
\setlength{\emergencystretch}{3em}
\newcommand{\N}{\mathbb{N}}
\newcommand{\Z}{\mathbb{Z}}
\newcommand{\Q}{\mathbb{Q}}
\newcommand{\R}{\mathbb{R}}
\newcommand{\C}{\mathbb{C}}
\newcommand{\F}{\mathbb{F}}
\newcommand{\A}{\mathbb{A}}
\newcommand{\PP}{\mathbb{P}}
\newcommand{\E}{\mathbb{E}}
\newcommand{\eps}{\varepsilon}
\newcommand{\dd}{\,\mathrm{d}}
\begin{document}
\section*{UnsolvedMath ID: 30002747}
\textbf{Source problem code:} catalogue entry
\subsection{Definitions and mathematical statement}
Let \(U\) denote the history-based mathematics unit under consideration, and let \(P\) denote its intended student population. These symbols name objects mentioned in the supplied record; the record does not specify a particular unit, population, or operational definition of ``history-based.'' An intervention here refers to implementing the unit with the intended students; its implementation conditions are unspecified.

The supplied statement asks for the following research program.
\begin{enumerate}
\item Determine how the ideas in \(U\) should be sequenced and which ways of thinking the unit affords.
\item Determine the intended students' relevant background and cognitive abilities, the feasible instructional time, and the unit's compatibility with current educational programs.
\item Under the intervention, determine students' learning trajectories, acquisition, retention, and later achievement.
\end{enumerate}

This is an empirical curriculum and learning research agenda. The supplied record does not give a mathematical proposition with specified domains and quantifiers. In particular, it specifies neither a criterion for selecting a sequence of ideas nor operational measures of ways of thinking, cognitive abilities, learning trajectories, acquisition, retention, or later achievement. It also leaves the instructional setting, observation times, and any comparison conditions unspecified. Consequently, a uniquely determined mathematical theorem, optimization problem, or statistical estimand cannot be extracted without additional assumptions. No universal claim over units or student populations, and no causal interpretation of outcomes under the intervention, is imposed here.
\subsection{Short English statement}
For a history-based mathematics unit, investigate its instructional design, suitability for the intended students and current programs, and students' learning trajectories, acquisition, retention, and later achievement during and after implementation. The source presents a research agenda whose population, measures, timing, and evaluation criteria require specification before a precise mathematical or statistical formulation is possible.
\subsection{Sources}
\begin{itemize}
\item Local registry: \texttt{data/unsolvedmath/problems.json}, permanent UnsolvedMath ID \(30002747\), source problem code \texttt{catalogue entry}, title ``Learning Trajectories in History-Based Mathematics Curricula.'' The supplied record assigns proposed year 2014 and status ``open.'' No individual proposer is supplied. Its statement status is ``exact,'' with the supplied verification note that the statement was internally readable, complete, and retained byte-for-byte; this is not independent verification against the report.
\item Supplied report citation: \emph{Mathematics in Undergraduate Study Programs: Challenges for Research and for the Dialogue between Mathematics and Didactics of Mathematics} (2015), Oberwolfach Reports, pp.\ 3103--3175, DOI 10.4171/owr/2014/56. Original OWR report: \url{https://doi.org/10.4171/owr/2014/56}. The supplied background describes Oberwolfach Reports as documenting mathematical workshops and their open-problem discussions.
\item Supplied literature reference: ``Kempen, design-based study of proof learning,'' \url{https://d-nb.info/1155236289/34}. The record provides no fuller bibliographic citation for this reference.
\item Literature assessment checked on \texttt{2026-08-21}: the supplied curation labels the agenda ``open'' and reports that the curriculum, population, retention, and achievement questions remain a research agenda. It states that a design-based study reports local course findings but does not settle the proposed learning trajectories or transfer effects in general. This is supplied dated curation, not independently verified literature assessment. The supplied background explicitly requests independent verification of time-sensitive status and cited claims.
\end{itemize}
\end{document}
